\begin{abstract}
    We propose a new model for the forecasting of both the implied volatility surfaces and the underlying asset price.
    In the spirit of \cite{guyon2022volatility} who are interested in the dependence of volatility indices (e.g. the VIX) on the paths of the associated equity indices (e.g. the S\&P 500), we first study how vanilla options implied volatility can be predicted using the past trajectory of the underlying asset price. Our empirical study reveals that a large part of the movements of the at-the-money-forward implied volatility for up to two years time-to-maturities can be explained using the past returns and their squares. Moreover, we show that this feedback effect gets weaker when the time-to-maturity increases. Building on this new stylized fact, we fit to historical data a parsimonious version of the SSVI parameterization (\citeauthor{gatheral2014arbitrage}, \citeyear{gatheral2014arbitrage}) of the implied volatility surface relying on only four parameters and show that the two parameters ruling the at-the-money-forward implied volatility as a function of the time-to-maturity exhibit a path-dependent behavior with respect to the underlying asset price. Finally, we propose a model for the joint dynamics of the implied volatility surface and the underlying asset price. The latter is modelled using a variant of the path-dependent volatility model of Guyon and Lekeufack and the former is obtained by adding a feedback effect of the underlying asset price onto the two parameters ruling the at-the-money-forward implied volatility in the parsimonious SSVI parameterization and by specifying Ornstein-Uhlenbeck processes for the residuals of these two parameters and Jacobi processes for the two other parameters. Thanks to this model, we are able to simulate highly realistic paths of implied volatility surfaces that are free from static arbitrage.
\end{abstract}

\keywords{Implied volatility modelling; SSVI; Path-dependent volatility; Simulation; Arbitrage}